%=========== ΛΥΣΗ - 5ο ΔΙΑΓΩΝΙΣΜΑ ΠΡΟΣΟΜΟΙΩΣΗΣ ===========
\begin{thema}{Α}
\begin{erwthma}
\item Θεωρούμε τη συνάρτηση
\[
h(x)=f(x)-g(x),\quad x\in\varDelta.
\]
Η $h$ είναι συνεχής στο $\varDelta$ και παραγωγίσιμη σε κάθε εσωτερικό σημείο του $\varDelta$. Επιπλέον
\[
h'(x)=f'(x)-g'(x)=0
\]
σε κάθε εσωτερικό σημείο του $\varDelta$.

Έστω $x_1,x_2\in\varDelta$ με $x_1<x_2$. Η $h$ είναι συνεχής στο $[x_1,x_2]$ και παραγωγίσιμη στο $(x_1,x_2)$, οπότε από το Θεώρημα Μέσης Τιμής υπάρχει $\xi\in(x_1,x_2)$ τέτοιο ώστε
\[
\frac{h(x_2)-h(x_1)}{x_2-x_1}=h'(\xi)=0.
\]
Άρα
\[
h(x_2)=h(x_1).
\]
Επομένως η $h$ είναι σταθερή στο $\varDelta$, δηλαδή υπάρχει $c\in\mathbb{R}$ τέτοιο ώστε
\[
h(x)=c
\]
για κάθε $x\in\varDelta$. Άρα
\[
f(x)-g(x)=c
\Leftrightarrow
f(x)=g(x)+c
\]
για κάθε $x\in\varDelta$.

\item Μια συνάρτηση $f$ παρουσιάζει τοπικό ελάχιστο στο σημείο $x_0$ του πεδίου ορισμού της, όταν υπάρχει περιοχή του $x_0$ τέτοια ώστε
\[
f(x)\geq f(x_0)
\]
για κάθε $x$ της περιοχής που ανήκει στο πεδίο ορισμού της $f$.

\item Η ευθεία $x=x_0$ λέγεται κατακόρυφη ασύμπτωτη της γραφικής παράστασης της $f$, όταν ένα τουλάχιστον από τα όρια
\[
\lim_{x\to x_0^-}f(x),\quad
\lim_{x\to x_0^+}f(x)
\]
είναι $+\infty$ ή $-\infty$.

\item Οι χαρακτηρισμοί είναι:
\begin{alist}
\item \textbf{Λάθος}. Η $f(x)=e^x$ είναι $1-1$, αλλά δεν έχει σημείο καμπής.
\item \textbf{Σωστό}. Για συνεχείς συναρτήσεις ισχύει
\[
\dint_a^b(f(x)+g(x))\,\d x=\dint_a^b f(x)\,\d x+\dint_a^b g(x)\,\d x.
\]
\item \textbf{Λάθος}. Για παράδειγμα, η $f(x)=x$ στο $(0,1)$ δεν έχει απόλυτο μέγιστο ούτε απόλυτο ελάχιστο.
\item \textbf{Λάθος}. Για παράδειγμα, η $f(x)=x^3$ έχει $f'(0)=0$, αλλά στο $x=0$ δεν παρουσιάζει τοπικό ακρότατο.
\item \textbf{Σωστό}. Για κάθε συνάρτηση για την οποία ορίζεται το ολοκλήρωμα ισχύει
\[
\dint_a^a f(x)\,\d x=0.
\]
\end{alist}
\end{erwthma}
\end{thema}

\begin{thema}{Β}
Δίνεται η συνάρτηση
\[
f(x)=x^3+x-2,\quad x\in\mathbb{R}.
\]
\begin{erwthma}
\item Για κάθε $x\in\mathbb{R}$ είναι
\[
f'(x)=3x^2+1>0.
\]
Άρα η $f$ είναι \textbf{γνησίως αύξουσα} στο $\mathbb{R}$.

Επιπλέον
\[
\lim_{x\to-\infty}f(x)=-\infty
\quad\text{και}\quad
\lim_{x\to+\infty}f(x)=+\infty.
\]
Επειδή η $f$ είναι συνεχής και γνησίως αύξουσα στο $\mathbb{R}$, έχει πεδίο τιμών
\[
{f(\mathbb{R})=\mathbb{R}}.
\]

Αφού $0\in f(\mathbb{R})$, η εξίσωση $f(x)=0$ έχει τουλάχιστον μία ρίζα. Επειδή η $f$ είναι γνησίως αύξουσα, η ρίζα είναι μοναδική. Μάλιστα
\[
f(1)=1+1-2=0.
\]
Άρα η $f$ έχει \textbf{ακριβώς μία ρίζα}, το $x=1$.

\item Επειδή η $f$ είναι γνησίως αύξουσα στο $\mathbb{R}$, είναι συνάρτηση $1-1$, άρα αντιστρέφεται με καθολικό τρόπο. Το πεδίο ορισμού της $f^{-1}$ είναι το σύνολο τιμών της $f$, δηλαδή
\[
D_{f^{-1}}=\mathbb{R}.
\]

Για το $f^{-1}(-2)$ έχουμε
\[
f^{-1}(-2)=y
\Leftrightarrow f(y)=-2
\Leftrightarrow y^3+y-2=-2
\Leftrightarrow y^3+y=0.
\]
Άρα
\[
y(y^2+1)=0
\Leftrightarrow y=0.
\]
Επομένως
\[
{f^{-1}(-2)=0}.
\]

Για το $f^{-1}(8)$ έχουμε
\[
f(2)=2^3+2-2=8.
\]
Επειδή η $f$ είναι $1-1$, προκύπτει
\[
{f^{-1}(8)=2}.
\]

\item Θέτουμε
\[
t=f^{-1}(x).
\]
Τότε
\[
x=f(t)=t^3+t-2.
\]
Όταν $x\to-2$, έχουμε
\[
t\to f^{-1}(-2)=0.
\]
Άρα
\begin{align*}
\lim_{x\to-2}\frac{f^{-1}(x)-0}{x+2}
&=\lim_{t\to0}\frac{t}{f(t)+2}\\
&=\lim_{t\to0}\frac{t}{t^3+t}\\
&=\lim_{t\to0}\frac{1}{t^2+1}=1.
\end{align*}
Επομένως
\[
{\lim_{x\to-2}\frac{f^{-1}(x)-0}{x+2}=1}.
\]

\item Στο διάστημα $[-2,8]$ η $f^{-1}$ είναι μη αρνητική, αφού είναι γνησίως αύξουσα και
\[
f^{-1}(-2)=0.
\]
Άρα το ζητούμενο εμβαδόν είναι
\[
E=\dint_{-2}^{8} f^{-1}(x)\,\d x.
\]
Θέτουμε
\[
x=f(t)=t^3+t-2.
\]
Τότε
\[
\d x=f'(t)\,\d t=(3t^2+1)\,\d t.
\]
Τα άκρα γίνονται:
\[
x=-2\Rightarrow t=0
\quad\text{και}\quad
x=8\Rightarrow t=2.
\]
Άρα
\begin{align*}
E&=\dint_0^2 t(3t^2+1)\,\d t\\
&=\dint_0^2 (3t^3+t)\,\d t\\
&=\left[\frac{3t^4}{4}+\frac{t^2}{2}\right]_0^2\\
&=12+2=14.
\end{align*}
Επομένως
\[
{E=14}.
\]
\end{erwthma}
\end{thema}

\begin{thema}{Γ}
Θεωρούμε τη συνάρτηση
\[
f(x)=(x-1)\ln x,\quad x>0.
\]
\begin{erwthma}
\item Για κάθε $x>0$ είναι
\[
f'(x)=\ln x+\frac{x-1}{x}
=\ln x+1-\frac{1}{x}.
\]
Θεωρούμε τη συνάρτηση
\[
h(x)=\ln x+1-\frac{1}{x},\quad x>0.
\]
Έχουμε
\[
h'(x)=\frac{1}{x}+\frac{1}{x^2}>0,
\]
άρα η $h$ είναι γνησίως αύξουσα στο $(0,+\infty)$. Επειδή
\[
h(1)=0,
\]
παίρνουμε
\[
h(x)<0\quad\text{για }0<x<1
\]
και
\[
h(x)>0\quad\text{για }x>1.
\]
Επομένως
\[
f'(x)<0\quad\text{για }0<x<1
\]
και
\[
f'(x)>0\quad\text{για }x>1.
\]

Στον παρακάτω πίνακα βλέπουμε τα πρόσημα της $f'$ και τη μονοτονία της $f$:
\begin{center}
\begin{tikzpicture}
\tkzTabInit[color,lgt=2,espcl=1.5,colorC = maincolor!50!white,colorL = maincolor!20!white,colorV = maincolor!50!white]%
{$x$ /0.8,$f'(x)$ /0.8,$f(x)$/0.8}%
{$0$,$1$,$+\infty$}%
\tkzTabLine{d, -, z, +, }
\tkzTabLine{d, \searrow,t,\nearrow,}
\end{tikzpicture}
\end{center}
Άρα η $f$ είναι \textbf{γνησίως φθίνουσα} στο $(0,1]$ και \textbf{γνησίως αύξουσα} στο $[1,+\infty)$.

Επιπλέον
\[
f(1)=0,
\]
οπότε η $f$ παρουσιάζει ολικό ελάχιστο στο $x_0=1$, με ελάχιστη τιμή $0$.

Για την κυρτότητα έχουμε
\[
f''(x)=\left(\ln x+1-\frac{1}{x}\right)'
=\frac{1}{x}+\frac{1}{x^2}
=\frac{x+1}{x^2}>0.
\]
Άρα η $f$ είναι \textbf{κυρτή} στο $(0,+\infty)$.

\item Από το προηγούμενο ερώτημα, η $f$ παρουσιάζει ολικό ελάχιστο στο $x=1$ με τιμή $0$. Άρα, για κάθε $x>0$,
\[
f(x)\geq f(1)=0.
\]
Επομένως
\[
{(x-1)\ln x\geq0,\quad x>0}.
\]
Η ισότητα ισχύει μόνο για
\[
{x=1}.
\]

\item Από το προηγούμενο ερώτημα, στο διάστημα $[1,e]$ ισχύει
\[
f(x)\geq0.
\]
Άρα το ζητούμενο εμβαδόν είναι
\[
E=\dint_1^e (x-1)\ln x\,\d x.
\]
Με ολοκλήρωση κατά παράγοντες, παίρνουμε
\begin{align*}
E&=\dint_1^e \left(\frac{x^2}{2}-x\right)'\ln x\,\d x\\
&=\left[\left(\frac{x^2}{2}-x\right)\ln x\right]_1^e
-\dint_1^e\left(\frac{x^2}{2}-x\right)\frac{1}{x}\,\d x\\
&=\frac{e^2}{2}-e-\dint_1^e\left(\frac{x}{2}-1\right)\,\d x\\
&=\frac{e^2}{2}-e-\left[\frac{x^2}{4}-x\right]_1^e\\
&=\frac{e^2}{2}-e-\left(\frac{e^2}{4}-e-\frac{1}{4}+1\right)\\
&=\frac{e^2}{4}-\frac{3}{4}.
\end{align*}
Άρα
\[
{E=\frac{e^2-3}{4}}.
\]

\item Η ευθεία
\[
y=\left(2-\frac{1}{e}\right)x-2026
\]
έχει συντελεστή διεύθυνσης
\[
\lambda=2-\frac{1}{e}.
\]
Ζητούμε σημείο της $C_f$ με τετμημένη $x_0>1$ τέτοιο ώστε
\[
f'(x_0)=2-\frac{1}{e}.
\]
Επειδή
\[
f'(x)=\ln x+1-\frac{1}{x},
\]
δοκιμάζουμε $x_0=e$ και βρίσκουμε
\[
f'(e)=1+1-\frac{1}{e}=2-\frac{1}{e}.
\]
Άρα η εφαπτομένη είναι παράλληλη στη δοσμένη ευθεία στο σημείο με τετμημένη $e$.

Η τεταγμένη του σημείου είναι
\[
f(e)=(e-1)\ln e=e-1.
\]
Επομένως το ζητούμενο σημείο είναι
\[
{M(e,e-1)}.
\]
\end{erwthma}
\end{thema}

\begin{thema}{Δ}
Θεωρούμε συνάρτηση $f:[0,2]\to\mathbb{R}$ με συνεχή δεύτερη παράγωγο, για την οποία
\[
f(0)=f'(0)=f'(2)=0
\]
και
\[
\dint_0^2 (x-2)f''(x)\,\d x=1.
\]
\begin{erwthma}
\item Με ολοκλήρωση κατά παράγοντες έχουμε
\begin{align*}
\dint_0^2 (x-2)f''(x)\,\d x
&=\left[(x-2)f'(x)\right]_0^2-\dint_0^2 f'(x)\,\d x\\
&=0-\bigl[f(x)\bigr]_0^2\\
&=-(f(2)-f(0))\\
&=-f(2).
\end{align*}
Επειδή το ολοκλήρωμα ισούται με $1$, προκύπτει
\[
-f(2)=1
\Leftrightarrow
{f(2)=-1}.
\]

\item Από το Θεώρημα Μέσης Τιμής για την $f$ στο $[0,2]$, υπάρχει $c\in(0,2)$ τέτοιο ώστε
\[
f'(c)=\frac{f(2)-f(0)}{2-0}
=\frac{-1-0}{2}
=-\frac{1}{2}.
\]
Θεωρούμε τη συνάρτηση
\[
\Phi(x)=3f'(x)+\frac{1}{1+f^2(x)}.
\]
Η $\Phi$ είναι συνεχής στο $[0,c]$. Επιπλέον
\[
\Phi(0)=3f'(0)+\frac{1}{1+f^2(0)}=1>0.
\]
Ακόμη
\[
\Phi(c)=3\left(-\frac{1}{2}\right)+\frac{1}{1+f^2(c)}
\leq -\frac{3}{2}+1
=-\frac{1}{2}<0.
\]
Άρα, από το θεώρημα \eng{Bolzano}, υπάρχει $\xi\in(0,c)\subset(0,2)$ τέτοιο ώστε
\[
\Phi(\xi)=0.
\]
Δηλαδή
\[
3f'(\xi)+\frac{1}{1+f^2(\xi)}=0
\Leftrightarrow
{3f'(\xi)=\frac{-1}{1+f^2(\xi)}}.
\]

\item Από το προηγούμενο ερώτημα του Θεωρήματος Μέσης Τιμής υπάρχει $c\in(0,2)$ τέτοιο ώστε
\[
f'(c)=-\frac{1}{2}.
\]
Εφαρμόζοντας το Θεώρημα Μέσης Τιμής για την $f'$ στο $[0,c]$, υπάρχει $u\in(0,c)$ τέτοιο ώστε
\[
f''(u)=\frac{f'(c)-f'(0)}{c-0}
=\frac{-\frac{1}{2}-0}{c}<0.
\]
Εφαρμόζοντας το Θεώρημα Μέσης Τιμής για την $f'$ στο $[c,2]$, υπάρχει $v\in(c,2)$ τέτοιο ώστε
\[
f''(v)=\frac{f'(2)-f'(c)}{2-c}
=\frac{0+\frac{1}{2}}{2-c}>0.
\]
Άρα η συνεχής συνάρτηση $f''$ παίρνει και αρνητικές και θετικές τιμές στο $(0,2)$. Επομένως η κυρτότητα της $f$ αλλάζει τουλάχιστον μία φορά στο $(0,2)$, άρα η γραφική παράσταση της $f$ παρουσιάζει τουλάχιστον ένα σημείο καμπής στο διάστημα $(0,2)$.

\item Θέτουμε
\[
\psi(x)=f'(x)+\frac{1}{2},\quad x\in[0,2].
\]
Η $\psi$ είναι συνεχής στο $[0,2]$ και
\[
\psi(0)=\psi(2)=\frac{1}{2}>0.
\]
Επίσης
\[
\dint_0^2\psi(x)\,\d x
=\dint_0^2 f'(x)\,\d x+\dint_0^2\frac{1}{2}\,\d x
=f(2)-f(0)+1
=-1+1=0.
\]
Αν ήταν $\psi(x)\geq0$ για κάθε $x\in[0,2]$, τότε, επειδή $\psi(0)>0$, θα είχαμε
\[
\dint_0^2\psi(x)\,\d x>0,
\]
άτοπο. Άρα υπάρχει $c\in(0,2)$ τέτοιο ώστε
\[
\psi(c)<0.
\]
Επειδή $\psi(0)>0$ και $\psi(c)<0$, από το θεώρημα \eng{Bolzano} υπάρχει $x_1\in(0,c)$ τέτοιο ώστε
\[
\psi(x_1)=0
\Leftrightarrow f'(x_1)=-\frac{1}{2}.
\]
Επειδή $\psi(c)<0$ και $\psi(2)>0$, υπάρχει $x_2\in(c,2)$ τέτοιο ώστε
\[
\psi(x_2)=0
\Leftrightarrow f'(x_2)=-\frac{1}{2}.
\]
Έτσι $x_1<x_2$ και
\[
\frac{1}{f'(x_1)}+\frac{1}{f'(x_2)}
=\frac{1}{-\frac{1}{2}}+\frac{1}{-\frac{1}{2}}
=-2-2=-4.
\]

\item Δίνεται ότι
\[
\lim_{x\to0}\frac{f(x)}{x^2}=-2.
\]
Έχουμε
\[
\lim_{x\to0}\frac{\hm f(x)}{f(x)}=1
\]
αφού $f(x)\to0$. Επίσης
\[
x^2f(x)+f^2(x)=f(x)(x^2+f(x)).
\]
Άρα
\begin{align*}
\frac{\hm f(x)}{x^2f(x)+f^2(x)}
&=\frac{\hm f(x)}{f(x)}\cdot\frac{1}{x^2+f(x)}.
\end{align*}
Όμως
\[
\frac{x^2+f(x)}{x^2}
=1+\frac{f(x)}{x^2}\to 1-2=-1.
\]
Άρα
\[
x^2+f(x)\sim -x^2
\]
και επομένως
\[
\frac{1}{x^2+f(x)}\to-\infty.
\]
Συνεπώς
\[
{\lim_{x\to0}\frac{\hm f(x)}{x^2f(x)+f^2(x)}=-\infty}.
\]
\end{erwthma}
\end{thema}
%=========== ΤΕΛΟΣ ΛΥΣΗΣ - 5ο ΔΙΑΓΩΝΙΣΜΑ ΠΡΟΣΟΜΟΙΩΣΗΣ ===========
